

\section{The index number problem and why we assume homogenous industry-level output}
\label{appendix:index_numbers}

To see why we prefer to assume homogenous output than to use index numbers, consider a simple but pathological example (Table \ref{tab_indices_example}). Our economy features a transition from a dominant dirty sector to a prominent clean sector, and the clean sector is initially more expensive but is getting cheaper. 

\input{tables/tab_indices_examples_p_q}

Because both prices are non-increasing, prices indices are non-increasing, even though the price of an ``average unit'' is increasing.

More precisely, with the numbers from Table \ref{tab_indices_example}, the Paasche price index would be $(50\times600+50\times 1200)/(50\times600+50\times1500)=0.86$, indicating a \textit{decrease} of roughly 15\% of the price, even though we have replaced 40\% of the conventional production by the expensive clean good. A decrease in price would translate, in our framework, into an increase in productivity. We consider this a pathological outcome given that we have replaced a cheap technology by an expensive one. 

We have evaluated in Table \ref{tab_indices_example} various price and quantity indices and the basic issue remains, so we prefer to assume that the dirty incumbent and the alternative produce a homogenous product. This is known as the ``unit value'' index in the index number literature \citep{balk2012price}. This index assumes homogeneity between the units, so that the price index is the average price paid for one unit of good, and the quantity index is the total number of units.

\input{tables/tab_indices_example}



